\documentclass[11pt]{amsart}

\usepackage[nomath,nothms,nomisc,nolayout,nohyperref,abbrs,nocolor]{mymacros}

\usepackage{amsmath,amssymb,bbm,mathtools}
\usepackage{xspace}
\usepackage{thmtools, thm-restate}
\usepackage{fullpage}
\usepackage{enumerate}
\usepackage{microtype}

\usepackage[bookmarks, colorlinks, breaklinks,
pdftitle={},pdfauthor={}]{hyperref}

\hypersetup{
  linkcolor=black,
  citecolor=black,
  filecolor=black,
  urlcolor=black}

\numberwithin{equation}{section}

\usepackage{cleveref}

%% Macros of h02rg.tex
\newcommand{\conn}{\leftrightarrow}
\newcommand{\nconn}{\not\leftrightarrow}
\newcommand{\obs}{\star}
\newcommand{\cVbulk}{\cV^\varnothing}
\newcommand{\cKbulk}{\cK^\varnothing}
\newcommand{\Vbulk}{V^\varnothing}
\newcommand{\Vobs}{V^\obs}
\newcommand{\uobs}{u^\obs}
\newcommand{\Kbulk}{K^\varnothing}
\newcommand{\Kobs}{K^\obs}
\newcommand{\ellabN}{\ell_{ab,N}}
\newcommand{\OL}{O_L}
\newcommand{\dplus}{(\frac{d-2}{2} \wedge 1)}
\newcommand{\Loc}{\operatorname{Loc}}
\newcommand{\loc}{\operatorname{loc}}
\newcommand{\tu}{\tilde u_{N,N}}
\newcommand{\Comp}{\operatorname{Comp}}

\title{Erratum: Percolation transition for random forests in $d\geq 3$}
\date{}

\author{Roland Bauerschmidt}
\address{University of Cambridge, DPMMS; Courant Institute of Mathematical Sciences, NYU}
\email{bauerschmidt@cims.nyu.edu}

\author{Nicholas Crawford}
\address{The Technion, Department of Mathematics}
\email{nickc@technion.ac.il}

\author{Tyler Helmuth}
\address{Durham University, Department of Mathematical Sciences}
\email{tyler.helmuth@durham.ac.uk}

\begin{document}
\sloppy

\begin{abstract}
  We correct a number of errors in \cite{BCH24}. None of them affects the main results,
  Theorems~1.1, 1.2, 1.3 and~1.5 and Proposition~1.4 of \cite{BCH24}.
\end{abstract}

\maketitle

\section{Introduction}

This note lists corrections to \cite{BCH24}. They were found during a formal verification of
the main results of \cite{BCH24} in the Lean proof assistant. Section, theorem and equation
numbers refer to \cite{BCH24}, and the notation is that of \cite{BCH24}.

The main results, Theorems~1.1, 1.2, 1.3, 1.5 and Proposition~1.4, hold as stated. Some
intermediate statements in Sections~2--6 need to be changed. Each correction below is marked
\textbf{(S)} if it changes the statement of a numbered result (or a definition), and
\textbf{(P)} if it only concerns a proof. Where a statement changes, we indicate why the main
results are not affected.

\section{Corrections}

\subsection*{Section 2: Consequences of combining renormalisation and Ward identities}

\begin{itemize}
  % h02rg.tex l. 1066-1073, thm:psi4, e:psi4-00
\item \textbf{(S)} Theorem~2.1, (2.8). The error term $O(b_0L^{-\kappa N})$ is false: for
  $b_0=0$ the left-hand side is the torus Green function $(-\Delta+m^2)^{-1}(0,0)$, which differs
  from $\gamma=(-\Delta^{\Z^d}+m^2)^{-1}(0,0)$ by the zero mode $t_N/|\Lambda_N|$ and by the
  difference between the torus and $\Z^d$ Green functions, both of order $L^{-(d-2)N}$ when
  $L^{-N}\le m$ (see also the correction of (6.26) below). The correct statement is
  \begin{equation*}
    \avg{\bar\psi_0\psi_0}_{m^2,s_0,a_0,b_0} = \gamma + O(L^{-(d-2)N}) + O(b_0L^{-\kappa N}).
  \end{equation*}
  Only the limit $N\to\infty$ of (2.8) is used (in the proof of Corollary~2.2(i)), so nothing
  else changes.

  % h02rg.tex l. 1103-1118, e:betab0m, e:hb0m, cor:ward
\item \textbf{(P)} Section~2.1, (2.14)--(2.15) and Corollary~2.2. It is used, but not shown,
  that $s_0^c(b_0,m^2)$ and $a_0^c(b_0,m^2)$ are real for real $b_0$; otherwise $\beta(b_0,m^2)$
  and $h(b_0,m^2)$ are not real parameters. This follows since the covariances are real, so
  the renormalisation group map commutes with complex conjugation of all coefficients, and
  hence, by uniqueness of the solution of the stable manifold problem in the proof of
  Theorem~3.24, $V_0^c(m^2,\bar b_0)=\overline{V_0^c(m^2,b_0)}$. The same argument shows that
  $\tilde a^c_{N,N}$, $\tilde u^c_{N,N}$, $u_N^c$, and $\lambda$, $\gamma$ at $m^2=0$, are real for
  real $b_0$, which is also used in Section~2.2.

  % h02rg.tex l. 1383-1391, cor:macro, proof of the second claim
\item \textbf{(P)} Proof of Lemma~2.4, second claim. That $1+\tu^c$ diverges as $m^2\downarrow 0$
  only gives $\tilde a^c_{N,N}(b_0,0)\ge 0$. Strict positivity follows from the rate:
  by (2.28) and $m^2t_N\to1$, $m^2(1+\tu^c)\to\tilde a^c_{N,N}(b_0,0)$ as $m^2\downarrow0$, and by
  (2.41) together with $\det(-\Delta+m^2)\le K_Nm^2$ (only one eigenvalue of $-\Delta+m^2$ tends
  to $0$),
  \begin{equation*}
    m^2(1+\tu^c) \geq K_N^{-1}\,Z_{\beta,h}\pa{\frac{\beta e^{-u_N^c}}{1+s_0^c}}^{-|\Lambda_N|},
  \end{equation*}
  whose right-hand side converges to a positive limit since $Z_{\beta,h}\to Z_{\beta,0}>0$.
\end{itemize}

\subsection*{Section 3: The bulk renormalisation group flow}

\begin{itemize}
  % h02rg.tex l. 1717-1719, sec:frd (zero mode); see also App. B
\item \textbf{(S)} Section~3.1, after (3.6). With the corrected definition of $C_N$ (see the
  correction to Appendix~B), $C_N$ satisfies (3.3) with $j+1=N$, i.e.,
  $|\nabla^\alpha C_N(x,y)|\le O_{\alpha,s}(\vartheta_{N-1}(m^2))L^{-(d-2+|\alpha|)(N-1)}$, not
  (3.3) with $j=N$. Since $\ell_0$ is chosen $L$-dependent in (3.24), this does not affect the
  rest of the paper.

  % h02rg.tex l. 2018-2021, def:Phi
\item \textbf{(S)} Definition~3.2. The test functions must be complex valued,
  $g\colon\Lambda_f^*\to\C$ (as in \cite{MR3332938}), since $\cN(\Lambda)$ has complex
  coefficients. With real test functions the product property (3.26) fails: for $\ell=1$ and
  $F=1+i\psi_x$, one gets $\|F\|_{T_j(\ell)}=\sqrt2$ but $\|F^2\|_{T_j(\ell)}=\|1+2i\psi_x\|_{T_j(\ell)}=\sqrt5>2$.

  % h02rg.tex l. 1958-1962 (def:cKbulk) and l. 2333-2361 (prop:Loc (v) and its proof)
\item \textbf{(S)} Proposition~3.8(v) is false as stated. For a block $B$, the element
  $F=\sum_{x\in B}\pa{\psi_x(\Delta\bar\psi)_x-(\Delta\psi)_x\bar\psi_x}\in\cN_{\rm sym}(B)$ is
  invariant under all automorphisms fixing $B$ and satisfies $\Loc_{B,B}F=F$, but
  $F\notin(\C\oplus\cVbulk)(B)$, since (3.35) contains only the symmetric combination of
  $(-\Delta\psi)\bar\psi$ and $\psi(-\Delta\bar\psi)$. Correspondingly, the claim in the proof that
  $\cVbulk$ contains all local polynomials of dimension at most $d$ that are invariant under
  the automorphisms fixing a point is false. The correction is to require in addition that
  the $F_i$ are invariant under the symplectic symmetry $J\colon\psi\mapsto\bar\psi$,
  $\bar\psi\mapsto-\psi$ (note $JF=-F$ for the above $F$). The maps $\Loc_{X,Y}$, $\theta$ and
  $\E_C$ (for symmetric $C$) commute with $J$, and the initial potential is $J$-invariant.
  Accordingly, $J$-invariance of each $K_j(X)$ is to be added to the definition of
  $\cKbulk_j$ (Definition~3.1); it is preserved by the renormalisation group map, so nothing
  else changes.

  % h02rg.tex l. 2560-2561, prop:ZVK, last sentence
\item \textbf{(S)} Proposition~3.11, last sentence. The automorphism invariance of $K_{j+1}$
  follows only for $j+1<N$: for $j+1=N$ every automorphism maps the single $N$-block to itself,
  but need not map $(N-1)$-blocks to $(N-1)$-blocks, so the invariance of $K_{N-1}$ does not
  give invariance of $K_N$ under all automorphisms. The invariance of $K_N$ is not used.

  % h02rg.tex l. 2818-2819, thm:step, last sentence
\item \textbf{(S)} Theorem~3.13, last sentence. The last map $\Phi_N$ satisfies the bounds with
  $L^{-\dplus}$ replaced by $L^d$, not by $1$: for $j+1=N$ there is no extraction ($Q=J=0$), and
  the $O(L^d)$ small sets $X$ all contribute $\E_{C_N}\theta K_{N-1}(X)$ to $K_N(\Lambda_N)$
  without contraction. Since $\Phi_N$ is applied only once, this is harmless.

  % h02rg.tex l. 3381-3388, lem:Ri-bd proof, e:K-crucial and the next line
\item \textbf{(P)} Proof of Lemma~3.22, (3.111) and the line after it. The factor
  $(1-e^{(u+V)_{j+1}(X_\cX)})$ should read $(1-e^{-(u+V)_{j+1}(X_\cX)})$.

  % h02rg.tex l. 3740-3743, proof of thm:flow, bound on D_j
\item \textbf{(P)} Proof of Theorem~3.24, the bound on $D_j$. The factor $b_0$ should be
  omitted: $\|D_j\|_{j\to j+1}\le\max\{L^{-(d-2)},O(L^{-3}+A^{-\eta})\}\le O(L^{-\kappa})$. This is
  the bound needed for the stable manifold theorem.
\end{itemize}

\subsection*{Section 4: Computation of the susceptibility}

\begin{itemize}
  % h02rg.tex l. 3906-3916, proof of prop:ZNN, e:kN2normbd
\item \textbf{(P)} Proof of Proposition~4.1, (4.10). With the pairing (3.25) and antisymmetric
  coefficients, the test function $g$ with $g_{(x,1),(y,-1)}=1$ and $g_z=0$ otherwise gives
  $\avg{K_N(\Lambda_N),g}=\frac12k_N^2$. The test function must be antisymmetrised,
  $g_{(x,1),(y,-1)}=1$ and $g_{(y,-1),(x,1)}=-1$; then $\avg{K_N(\Lambda_N),g}=k_N^2$ and still
  $\|g\|_{\Phi_N(\ell_N)}=\ell_N^{-2}$, so (4.10) holds.
\end{itemize}

\subsection*{Section 5: Additivity of the observable potential}

% h02rg.tex l. 4123-4160 (def:cVobs, e:VobsX); affects sec:locobs (prop:Loc-obs (iv)),
% sec:rgmap-ext l. 4672-4675 (prop:ZVK, prop:consistency for the extended map), lem:ZNN-obs
\noindent\textbf{(S)} Definition~5.1, (5.9). With the indicator $1_{a\in X,b\in X}$ on the
$\sigma_a\bar\sigma_b$ terms (and analogously in Case~(2)), the map
$X\mapsto(\uobs+\Vobs)(X)$ is not additive over blocks: if $a$ and $b$ lie in different blocks
$B_a\ne B_b$, then $(\uobs+\Vobs)(B_a\cup B_b)\ne(\uobs+\Vobs)(B_a)+(\uobs+\Vobs)(B_b)$. Additivity
is needed both in the polymer formalism, which uses $e^{-V(Y)}=\prod_{B\in\cB_j(Y)}e^{-V(B)}$
(proof of Proposition~3.11, extended in Section~5.5), and in Proposition~5.6(iv).

\emph{Corrected definition.} The $\sigma_a\bar\sigma_b$ terms (in Case~(2) the $\sigma_a\sigma_b$
terms) are evaluated block-additively: half of each such term is assigned to the block
containing $a$ and half to the block containing $b$. That is, in Case~(1),
\begin{multline*}
  (\uobs+\Vobs)(X) =
  -\lambda_a\sigma_a\bar\psi_a1_{a\in X}-\lambda_b\psi_b\bar\sigma_b1_{b\in X}
  \\
  +\sigma_a\bar\sigma_b\pa{-\frac q2+\frac r2\bar\psi_a\psi_a}1_{a\in X}
  +\sigma_a\bar\sigma_b\pa{-\frac q2+\frac r2\bar\psi_b\psi_b}1_{b\in X},
\end{multline*}
and analogously in Case~(2), where also the $\eta$ term is split in halves. This agrees with
(5.9) on every $X$ that contains both or neither of $a,b$, in particular on $X=\Lambda_N$, so
all statements about $\Lambda_N$ (e.g.\ (5.6) and Lemma~6.3) are unchanged. It entails the
following changes.
\begin{enumerate}
\item The $\sigma_a\bar\sigma_b$ part of the localisation of $K_j(X)$ on a small set $X$ containing
  exactly one of $a,b$, say $a$, is $\loc^{ab}_{X\cap\{a,b\},\cdot}$ with $X\cap\{a,b\}=\{a\}$; it
  produces a constant and a multiple of $\bar\psi_a\psi_a$ only, in the half at $a$. To accommodate
  this, $q$ and $r$ are split into $q_a,q_b$ and $r_a,r_b$ (the terms
  $\sigma_a\bar\sigma_b(-q_x+r_x\bar\psi_x\psi_x)$ in the half at $x\in\{a,b\}$), and in the
  recursions of $q$ in Lemmas~5.3 and~5.13 the term $-r_jC_{j+1}(0,0)$ (see the correction below)
  becomes $-(r_{a,j}C_{j+1}(a,a)+r_{b,j}C_{j+1}(b,b))$. Only $q=q_a+q_b$ enters on sets containing
  both points. If the $ab$ component of $K_j(X)$ vanishes on small sets containing only one
  of $a,b$, the split is harmless ($r_a=r_b=r/2$).
\item In Case~(2), the half of the $\eta$ term assigned to the block of $a$,
  $-\sigma_a\sigma_b\frac\eta4(\bar\psi_a\psi_b+\bar\psi_b\psi_a)$, is not local at scale $j$ if $a$
  and $b$ lie in different $j$-blocks. The locality $K_{j+1}(X)\in\cN(X^\square)$ of the output
  still holds if $\eta_j=0$ or $j\ge j_{ab}$ (then every $j$-block containing $a$ or $b$ has both
  points in its scale-$(j+1)$ neighbourhood), which covers the whole flow since $\eta_j=0$ for
  $j<j_{ab}$.
\item At $j=j_{ab}$, the factorisation of $K_{j+1}$ over connected components asserted for the
  extended map in Section~5.5 does not hold in general, since $e^{-(u+V)_{j+1}(B)}$ need not be
  local. It does hold along the flow, i.e., $K_{j+1}(U)=\prod_{Y\in\Comp(U)}K_{j+1}(Y)$ for the
  activities produced from $(V_j,K_j)$, which is what is needed to iterate Proposition~3.11
  up to scale $N$ in Lemma~6.3: $e^{-(u+V)_{j+1}}$ enters $K_{j+1}$ only through the inclusion
  $\cN(\Lambda)\subset\cN(\Lambda\sqcup\Lambda)$ and is not integrated against $C_{j+1}$, so its
  lack of locality does not affect the factorisation.
\end{enumerate}

\subsection*{Section 5: The observable renormalisation group flow}

\begin{itemize}
  % h02rg.tex l. 4181, 4193 (lem:Vobs-flow-free), l. 5285, 5297 (lem:Vobs-recursion), l. 5990 (proof of lem:2pt)
\item \textbf{(S)} Lemma~5.3 (both cases, including (5.15)), Lemma~5.13 (both cases) and the
  proof of Lemma~6.5. The $r$ term in the recursion of $q$ has the opposite sign: $+r_jC_{j+1}(0,0)$
  should read $-r_jC_{j+1}(0,0)$, as follows from Definition~5.1 and is consistent with the term
  $-r_Nt_N|\Lambda_N|^{-1}$ in (6.10) and (6.13). This affects no estimate.

  % h02rg.tex l. 5346-5350, lem:Vobs-bd, e:Vobs-bd6
\item \textbf{(S)} Lemma~5.14, (5.89). Since $\|V\|_j$ is the supremum of $\|V(B)\|_{T_j(\ell_j)}$
  over blocks and the norm (5.24) is the sum of the norms of the components, a block containing
  both $a$ and $b$ contributes $|\lambda_{a,j}|+|\lambda_{b,j}|$. The bound should read
  $\|\Vobs_j\|_j\le2\lambda_0+\OL(\lambda_0^2)+\OL(M\lambda_0b_0)$. This changes only constants.
\end{itemize}

\subsection*{Section 6: Computation of pointwise correlation functions}

\begin{itemize}
  % h02rg.tex l. 5599-5601 (e:ZNN-lambda-b) and l. 5694 (proof)
\item \textbf{(S)} Lemma~6.3, (6.9). Since $\psi\bar\sigma_b=-\bar\sigma_b\psi$, the coefficient is
  $\tilde Z^{\bar\sigma_b\psi}_{N,N}=-\lambda_{b,N}+\OL(\ell_{b,N}^{-1}\ell_N^{-1}\|\Kobs_N\|_N)$.

  % h02rg.tex l. 5603-5609 (e:ZNN-q-a) and l. 5692-5696 (proof, Case (1), B)
\item \textbf{(S)} Lemma~6.3, (6.10), and its proof in Case~(1). The expression for $B$ omits the
  constant term $k_N^{\sigma_a\bar\sigma_b}$ (the coefficient of $\sigma_a\bar\sigma_b$ without
  $\psi,\bar\psi$ in $\Kobs_N(\Lambda_N)$), which is included in Case~(2). Hence (6.10) holds
  with an additional error $\OL(\ellabN^{-1}\|\Kobs_N\|_N)$. Consequently (6.38) holds with an
  additional error $\OL(\frac{b_0}{\lambda_0}L^{-(d-2+\kappa)N})/(1+\tu)$, and (6.4) with
  $\OL(b_0L^{-(d-2+\kappa)N})/(1+\tu)$.

  % h02rg.tex l. 5784-5788 (e:gammainfbd), consequences: e:2ptaa (l. 5794), e:2ptaa-bis (l. 5514),
  % e:2pt-b (l. 5951), e:2pt-b-bis (l. 5544), e:1pt-b-thm / E_00 (l. 1263, 1288), e:2pt-b-thm / E_00xx (l. 1266, 1291)
\item \textbf{(S)} Lemma~6.4, (6.26). The bound on $\gamma_{x,j}$ is false: by the proof,
  $\gamma_{x,\infty}-\gamma_{x,j}$ contains $\lambda_{x,\infty}\sum_{k\ge j}C_{k+1}(0,0)$, which is of
  order $\lambda_0L^{-(d-2)j}$ and has no factor $b_0$ (already for $b_0=0$,
  $\gamma_{x,j}=\lambda_0\sum_{k<j}C_{k+1}(0,0)$). The correct bound is
  \begin{equation*}
    \gamma_{x,j}=\gamma_\infty+\OL(\lambda_0L^{-(d-2)j}).
  \end{equation*}
  Consequently, (6.27), (6.3) and (2.30) (i.e., $E_{00}$) hold only with an additional error term
  $\OL(L^{-(d-2)N})$ (not divided by $1+\tu$), and (6.39), (6.5) and (2.31) (i.e., $E_{00xx}$) only
  with an additional error term $\OL(L^{-(d-2)N})(1+(m^2|\Lambda_N|)^{-1}/|1+\tu|)$; alternatively,
  these statements hold as stated with $\gamma$ replaced by the scale-$N$ couplings
  $\gamma^{(2)}_{x,N}/\lambda_0$. The main results are not affected: in the proof of
  Corollary~2.2 the limit $N\to\infty$ is taken, and in Lemmas~2.5 and~2.6 (where $m^2\downarrow0$)
  the additional terms can be absorbed into the existing relative errors $O(b_0L^{-\kappa N})$,
  using $|\Lambda_N|\tilde a^c_{N,N}(b_0,0)=O(b_0)$ and $\kappa\le d-2$.

  % h02rg.tex l. 5900-5911, end of proof of lem:1pt (point independence of lambda)
\item \textbf{(S)} Proof of Lemma~6.4, last paragraph. The argument shows that
  $\lambda^{(2)}_{a,\infty}$ is independent of $a$ only at $m^2=0$, since it takes
  $m^2\downarrow0$ first. For $m^2>0$, the statements (2.10), (2.31), (6.5) and (6.39) are proved
  only with $\lambda^2$ replaced by $\lambda_a\lambda_b$ and $2\lambda\gamma$ by
  $(\lambda_a+\lambda_b)\gamma$ (with $a=0$, $b=x$ in (2.10) and (2.31)), where $\lambda_x=\lambda^{(2)}_{x,\infty}/\lambda_0=1+O(b_0)$. At $m^2=0$,
  which is the case needed in Section~2.2 and for the main results, $\lambda_a=\lambda_b$.

  % h02rg.tex l. 5951-5966 (e:2pt-b), l. 5544-5554 (e:2pt-b-bis), l. 1266-1272 / 1291-1295 (e:2pt-b-thm, E_00xx)
\item \textbf{(S)} Lemma~6.5, (6.39); Proposition~6.2, (6.5); Theorem~2.3, $E_{00xx}$ in (2.31).
  The error term $\OL(b_0L^{-(d-2+\kappa)N})$ (resp.\ $\OL(\frac{b_0}{\lambda_0}L^{-(d-2+\kappa)N})$)
  arises from the constants $k_N^{\#}$ in (6.13) and is therefore divided by $1+\tu$; it should
  read $\OL(b_0L^{-(d-2+\kappa)N})/(1+\tu)$. The two forms agree when $|\tu|$ is small, e.g.\ under
  the assumption $L^{-N}\le m$ of Theorem~2.1.
  % h02rg.tex l. 5977-5979 (e:lambdajab), l. 6000 and 6031 (|a-b| <= L), l. 6065 (||K^obs_N||)
\item \textbf{(P)} Proof of Lemma~6.5. In (6.40), $\OL(\lambda_0b_0L^{-\kappa})$ should read
  $\OL(\lambda_0b_0L^{-\kappa j})$ (twice); the two occurrences of ``$|a-b|\le L$'' should read
  $L^{j_{ab}}\le2|a-b|_\infty<L^{j_{ab}+1}$ (cf.\ (5.21)); and in Case~(2),
  $\|\Kobs_N\|_N\le\OL(b_0\lambda_0L^{-\kappa})$ should read $\|\Kobs_N\|_N\le\OL(b_0\lambda_0L^{-\kappa N})$.

\end{itemize}

\subsection*{Appendix B: Finite range decomposition}

\begin{itemize}
  % h02rg.tex l. 6496, 6511, 6513 and the estimates l. 6519-6538
\item \textbf{(S)} Appendix~B. The lower limit $\frac12L^N$ of the $t$-integrals defining
  $C_{N,N}$, $C_N$ and $t_N$ must be $\frac12L^{N-1}$, the upper limit of the integral defining
  $C_{N-1}$; otherwise the scales $t\in(\frac12L^{N-1},\frac12L^N]$ belong to no covariance and
  (3.1) fails. In the subsequent estimates $L^N$ is replaced by $L^{N-1}$; in particular
  $|C_N(x,y)|\lesssim L^{-(d-2)(N-1)}$ and $t_N=1/m^2-O(L^{2(N-1)})$, which is (3.7). See also the
  correction to Section~3.1.
\end{itemize}

\bibliographystyle{plain}
\bibliography{all}

\end{document}
